%% sections/03_related_work.tex

\section{관련 연구}
\label{sec:related}

\subsection{투기적 디코딩과 워크로드 의존성}

투기적 디코딩은 Stern~\emph{et al.}~\cite{stern2018blockwise} 이후
Leviathan~\emph{et al.}~\cite{leviathan2023spec},
Chen~\emph{et al.}~\cite{chen2023spec} 에 의해 형식화되었고,
Medusa~\cite{cai2024medusa}, EAGLE 계열~\cite{li2024eagle,li2025eagle3},
모델 없는 Prompt Lookup~\cite{saxena2024prompt} 과
SuffixDecoding~\cite{snowflake-arctic} 으로 발전했다.

프로덕션 보고는 이 가속이 워크로드$\cdot$동시성에 강하게 의존함을 일관되게 드러낸다.
Red Hat 의 vLLM EAGLE-3 분석~\cite{redhat2025eagle3vllm} 은 RAG$\cdot$math 에서 강한 가속,
번역에서 저하, 고배치에서 회귀를 보고하고,
Amazon$\cdot$NVIDIA 의 P-EAGLE~\cite{peagle2026vllm} 은 B200 에서 동시성이
$1\to64$ 로 증가하면 가속이 $1.55\times\to1.05\times$ 로 감소함을 보인다.
본 연구는 이 의존성을 \emph{언제 켜는가}의 예측 조건($\alpha$ 대 overhead)으로
정량화하고, 특히 \emph{추론 특화/MoE 모델}에서의 구조적 실패를
대규모로 특성화한다는 점에서 차별된다.

\subsection{워크로드 인지 라우팅}

투기적 방법 선택을 다루는 연구로
Dual-Pool Token-Budget Routing~\cite{liu2026dualpool},
arXiv 2505.08600~\cite{specrouting2025}, Nightjar~\cite{nightjar2025} 가 있다.
llm-d~\cite{llmd2025} 는 Kubernetes 기반 disaggregation 으로
prefill/decode 분리와 워크로드 라우팅을 지원하나 복잡한 인프라를 요구한다.
본 연구는 llm-d 를 동일 B200 환경에서 \emph{실측 베이스라인}으로 비교하여
(\S\ref{subsec:res-llmd}), 클러스터 라우팅의 추가 이득이
최선 단일 방법 대비 14\%(10/70)에 그침을 보인다.

\subsection{이종 CPU-GPU 서빙}

호스트 CPU 활용은 CPU-as-accelerator~\cite{xu2026arclight},
SmartNIC 기반 CPU-free~\cite{siavashi2026blink} 로 나뉘며,
Chung~\emph{et al.}~\cite{chung2026cpuslowdowns} 은 멀티-GPU 추론의 주요 병목이
CPU 임을 정량화했다. 본 연구는 spec-decode 가 GPU 를 un-saturate 시켜
시스템을 host-bound 로 전환함을 측정으로 보이고(\S\ref{subsec:res-resource}),
유휴 CPU 가 그 host 병목을 떠안는 reclamation 방향을 제시한다(\S\ref{sec:direction}).
이는 CPU 로 GPU 매트멀을 흉내내는 접근과 달리, spec-decode 가 \emph{만든}
host-path 작업을 CPU 로 이전한다는 점에서 구분된다.
